% ---------------------------------------------------------------
% Replaces "Learning rate scheduling" and "Domain switching optimization scheme".
% Domain switching is removed: it presupposes two networks with differing diffusion
% coefficients either side of the interface, which the Robin formulation does not have.
% ---------------------------------------------------------------

\subsubsection{Learning rate scheduling}

First order methods have longer time-to-solution but require less memory, while second
order methods converge in fewer iterations at a much larger memory footprint. In
JAX-DIPS we primarily use first order methods such as Adam~\cite{kingma2014}. The
GMRES and conjugate gradient methods traditionally applied to sparse linear systems sit
between the two, building basis vectors from mutually conjugate gradients,
$p_j^{T} A p_i = 0$, and converging in at most $n$ steps.

We use Adam with a cosine-annealed learning rate,
\begin{equation}
    r_k = \frac{r_0}{2}\left[1 + \cos\!\left(\frac{\pi k}{K}\right)\right],
    \label{eq:cosine}
\end{equation}
where $k$ is the optimisation step and $K$ the total number of steps in the run. We
take $r_0 = 10^{-3}$ and $K$ equal to the full training horizon, so that the schedule
reaches zero exactly at the end of training irrespective of the epoch budget. This
removes the decay-rate tuning required by an exponential schedule, which never reaches
zero and whose endpoint depends on the number of epochs. The schedule is flat at both
ends and steepest in the middle, providing full-rate exploration early and a long
low-rate refinement phase at the end.

\subsubsection{Gradient scale and the optimiser stability constant}
\label{sec:eps}

Preconditioning the residual has a consequence for the optimiser that, to our
knowledge, has not been reported, and which we found to control the accuracy attainable
at fine resolutions.

The assembled finite-volume row at an interior cell is $\mathcal{O}(\Delta x^{3})$,
while the Jacobi diagonal is $d_{ijk} \approx 6\mu\,\Delta x$. The preconditioned
residual corresponding to a solution error $e$ is therefore
$\mathcal{O}(\Delta x^{2}\, e)$ at interior cells and $\mathcal{O}(\Delta x\, e)$ at
cells the interface crosses. Interface cells dominate the loss, and their number scales
as $\Delta x^{-2}$ against $\Delta x^{-3}$ for the interior, so that the gradient of
\eqref{eq:loss} with respect to the network parameters scales as
\begin{equation}
    \left\| \nabla_{\theta}\mathcal{L} \right\| \;\sim\; C\,\Delta x^{3}\, e .
    \label{eq:gradscale}
\end{equation}
Refining the grid thus reduces the loss's sensitivity to a given solution error: the
objective becomes a progressively weaker measure of the quantity one is trying to
minimise.

Adam updates parameters by $\theta \leftarrow \theta - r_k\,\hat{m}/(\sqrt{\hat{v}} +
\varepsilon)$, in which $\varepsilon$ exists solely to prevent division by zero and is
conventionally set to $10^{-8}$. That value presupposes gradients of order unity. Under
\eqref{eq:gradscale} the gradients at $N_x = 64$ are of order $10^{-8}$, so
$\varepsilon$ ceases to be a guard and becomes the dominant term of the denominator.
Adam's scale invariance is lost and the step is throttled by a factor
$\sqrt{\hat{v}}/(\sqrt{\hat{v}} + \varepsilon)$. Since the gradient falls with the
error, the throttle tightens as training proceeds, and optimisation stalls at a
solution-error floor
\begin{equation}
    e_{\rm floor} \;\sim\; \frac{\varepsilon}{C\,\Delta x^{3}},
    \label{eq:floor}
\end{equation}
which \emph{grows} as the grid is refined. The floor is negligible at coarse
resolutions, where the truncation error dominates, and becomes controlling at fine
ones -- producing error that increases with resolution, a symptom easily mistaken for a
defect of the discretisation.

Setting $\varepsilon = 10^{-16}$, well below the smallest gradient the run produces,
restores scale invariance. The change is inert wherever the gradients are large: for
$|g| \gtrsim 10^{-4}$ the two values give numerically identical steps, and $\varepsilon$
can only ever reduce a step, never enlarge it beyond $r_k$. Its effect is therefore
confined to the regime in which the default is harmful. In our sweep the solution error
at $N_x = 64$ fell from $2.77 \times 10^{-2}$ to $1.48 \times 10^{-3}$, a factor of
$18.7$, with all other settings unchanged; the two coarsest resolutions, being
truncation-limited, were unaffected, as \eqref{eq:floor} predicts.

We note that the same consideration applies to any absolute tolerance applied to this
loss. A quasi-Newton solver terminating on a gradient-norm tolerance of $10^{-3}$, a
common default, satisfies its convergence test at the first iteration.
